%-------------------------
% Resume in Latex
% Author : Jake Gutierrez
% Based off of: https://github.com/sb2nov/resume
% License : MIT
%------------------------

\documentclass[letterpaper,11pt]{article}

\usepackage{latexsym}
\usepackage[empty]{fullpage}
\usepackage{titlesec}
\usepackage{marvosym}
\usepackage[usenames,dvipsnames]{color}
\usepackage{verbatim}
\usepackage{enumitem}
\usepackage[hidelinks]{hyperref}
\usepackage{fancyhdr}
\usepackage[english]{babel}
\usepackage{tabularx}
\input{glyphtounicode}


%----------FONT OPTIONS----------
% sans-serif
% \usepackage[sfdefault]{FiraSans}
% \usepackage[sfdefault]{roboto}
% \usepackage[sfdefault]{noto-sans}
% \usepackage[default]{sourcesanspro}

% serif
% \usepackage{CormorantGaramond}
% \usepackage{charter}


\pagestyle{fancy}
\fancyhf{} % clear all header and footer fields
\fancyfoot{}
\renewcommand{\headrulewidth}{0pt}
\renewcommand{\footrulewidth}{0pt}

% Adjust margins
\addtolength{\oddsidemargin}{-0.5in}
\addtolength{\evensidemargin}{-0.5in}
\addtolength{\textwidth}{1in}
\addtolength{\topmargin}{-.5in}
\addtolength{\textheight}{1.0in}

\urlstyle{same}

\raggedbottom
\raggedright
\setlength{\tabcolsep}{0in}

% Sections formatting
\titleformat{\section}{
  \vspace{-4pt}\scshape\raggedright\large
}{}{0em}{}[\color{black}\titlerule \vspace{-5pt}]

% Ensure that generate pdf is machine readable/ATS parsable
\pdfgentounicode=1

%-------------------------
% Custom commands
\newcommand{\resumeItem}[1]{
  \item\small{
    {#1 \vspace{-2pt}}
  }
}

\newcommand{\resumeSubheading}[4]{
  \vspace{-2pt}\item
    \begin{tabular*}{0.97\textwidth}[t]{l@{\extracolsep{\fill}}r}
      \textbf{#1} & #2 \\
      \textit{\small#3} & \textit{\small #4} \\
    \end{tabular*}\vspace{-7pt}
}

\newcommand{\resumeSubSubheading}[2]{
    \item
    \begin{tabular*}{0.97\textwidth}{l@{\extracolsep{\fill}}r}
      \textit{\small#1} & \textit{\small #2} \\
    \end{tabular*}\vspace{-7pt}
}

\newcommand{\resumeProjectHeading}[2]{
    \item
    \begin{tabular*}{0.97\textwidth}{l@{\extracolsep{\fill}}r}
      \small#1 & #2 \\
    \end{tabular*}\vspace{-7pt}
}

\newcommand{\resumeSubItem}[1]{\resumeItem{#1}\vspace{-4pt}}

\renewcommand\labelitemii{$\vcenter{\hbox{\tiny$\bullet$}}$}

\newcommand{\resumeSubHeadingListStart}{\begin{itemize}[leftmargin=0.15in, label={}]}
\newcommand{\resumeSubHeadingListEnd}{\end{itemize}}
\newcommand{\resumeItemListStart}{\begin{itemize}}
\newcommand{\resumeItemListEnd}{\end{itemize}\vspace{-5pt}}

%-------------------------------------------
%%%%%%  RESUME STARTS HERE  %%%%%%%%%%%%%%%%%%%%%%%%%%%%


\begin{document}

%----------HEADING----------
\begin{center}
    \textbf{\Huge \scshape Gokul Raj Santhosh} \\ \vspace{1pt}
    \small (+91) 7561867413 $|$ \underline{gokuljalaja@gmail.com} $|$
    \underline{LinkedIn: linkedin.com/in/gokul-raj-santhosh-14681b174} $|$
    \underline{GitHub: github.com/gokulbot}
\end{center}




%-----------SUMMARY-----------
\section{Summary}
Robotics Engineer with 3 years of experience in \textbf{Visual Inertial SLAM}, \textbf{sensor fusion}, and \textbf{real-time state estimation} for autonomous systems. Built production-grade \textbf{localization and perception pipelines} in \textbf{C++/ROS}, deployed on \textbf{NVIDIA Jetson Orin}, with hands-on experience in \textbf{IMU pre-integration}, \textbf{loop closure}, \textbf{pose graph optimization}, and \textbf{deep learning-based perception}. Passionate about pushing the boundaries of \textbf{autonomous navigation} in real-world environments.

%-----------EDUCATION-----------
\section{Education}
  \resumeSubHeadingListStart
    \item
    \begin{tabular*}{0.97\textwidth}[t]{l@{\extracolsep{\fill}}r}
      \textbf{Indian Institute of Technology, Madras} & 2023 \\
    \end{tabular*}\vspace{-7pt}
    \small
    \begin{itemize}[leftmargin=0in, label={}]
      \item Dual Degree in Engineering \hfill CGPA: 8.56
      \item BTech: Mechanical Engineering
      \item MTech: Robotics
    \end{itemize}
    \vspace{-5pt}
  \resumeSubHeadingListEnd


%-----------EXPERIENCE-----------
\section{Experience}
  \resumeSubHeadingListStart

    \resumeSubheading
      {Robotics Engineer}{Jan 2025 -- Present}
      {Addverb}{Noida, India}
      
    \vspace{2pt}
    \textit{\small \textbf{Simulation Pipeline}}
    \resumeItemListStart
      \resumeItem{Developed a comprehensive \textbf{simulation pipeline} in \textbf{Isaac Sim} for \textbf{Trakr} quadruped robot, integrating \textbf{stereo cameras, LiDAR, and IMU} sensor suite for validating \textbf{Navigation and Perception algorithms}}
      \resumeItem{Trained a \textbf{PPO-based locomotion controller} using \textbf{RSL RL} and \textbf{Legged Gym} within \textbf{Isaac Lab}, designing \textbf{custom reward functions} and tuning \textbf{hyperparameters} to achieve stable, robust locomotion}
    \resumeItem{\textbf{Open-sourced} the complete pipeline, enabling continuous algorithm development during hardware downtime and accelerating team collaboration}
    \resumeItemListEnd
      
\vspace{2pt}
    \textit{\small \textbf{Visual Inertial SLAM}}
    \resumeItemListStart
      \resumeItem{Implemented a \textbf{Visual Inertial SLAM pipeline} with tightly-coupled \textbf{IMU pre-integration}, \textbf{sliding window optimization}, and a \textbf{mapping module} for real-time localization and mapping}
      \resumeItem{Built a \textbf{loop closure module} with global descriptor retrieval (\textbf{CosPlace} + \textbf{FAISS}), geometric verification via \textbf{SuperPoint} + \textbf{LightGlue}, and pose graph optimization with \textbf{GNC outlier rejection}}
      \resumeItem{Implemented \textbf{map-based relocalization} using a persisted descriptor map, enabling robust \textbf{tracking recovery} after odometry failure}
      \resumeItem{Handled \textbf{dynamic objects} by integrating async \textbf{YOLOv8-seg} segmentation with confidence-based feature downweighting, improving robustness in human-present environments}
      \resumeItem{Deployed the complete software stack on \textbf{NVIDIA Jetson Orin}, achieving real-time operation at \textbf{20 Hz}; evaluated on \textbf{EuRoC Machine Hall} sequences with translational \textbf{RMSE of 0.28\%} and rotation error of \textbf{0.020\textdegree/m}}
      \resumeItem{As a \textbf{POC}, upgraded the visual front-end to \textbf{SuperPoint} + \textbf{LightGlue} (TensorRT), reducing translation RMSE by \textbf{21.9\%} and rotation RMSE by \textbf{9.5\%}}
    \resumeItemListEnd
          

   % \vspace{2pt}
   %    \textit{\small \textbf{Uncertainty-Aware Visual SLAM}$^{*}$}
   %    \resumeItemListStart
   %      \resumeItem{Implemented \textbf{deep learning-based visual odometry} inspired by \textbf{\href{https://mac-vo.github.io/}{\color{blue}MAC-VO}}, utilizing \textbf{learned feature extractors} with \textbf{keypoint covariance estimation} for improved pose optimization}
   %      \resumeItem{Integrated learned VO model into \textbf{SLAM framework}, leveraging \textbf{uncertainty-aware features} to enhance robustness in challenging scenarios including texture-less and dynamic environments}
   %    \resumeItemListEnd

   \resumeSubheading
      {Perception Engineer}{June 2023 -- Jan 2025}
      {Monarch Tractor}{Hyderabad, India}
      
      \vspace{2pt}
      \textit{\small \textbf{Visual SLAM}}
      \resumeItemListStart
        \resumeItem{Conducted an in-depth literature review on \textbf{Visual Inertial SLAM} techniques such as \textbf{ORB-SLAM3}, \textbf{S-PTAM}, \textbf{Open VINS} and \textbf{VINS-Fusion} and benchmarked the same for translation and rotation errors as well as computational efficiency}
        \resumeItem{Tuned and deployed these SLAM algorithms by writing \textbf{ROS Wrappers} and \textbf{docker containers} while achieving a translational \textbf{RMSE} of \textbf{2.7711\%} and average rotation error of \textbf{0.160 deg/m}}
      \resumeItemListEnd
      
      \vspace{2pt}
      \textit{\small \textbf{Inertial Odometry}}
      \resumeItemListStart
        \resumeItem{Improved GPS devoid localization by developing a \textbf{GRU}-based \textbf{DNN} tailored for inertial navigation}
        \resumeItem{Performed comprehensive data collection and cleaning to train the neural network, resulting in a mean \textbf{translation error} of \textbf{0.20 m} and a mean \textbf{rotational error} of \textbf{1.87 degrees} at the end of a \textbf{10 second sequence}}
        \resumeItem{Successfully deployed the model on a fleet of \textbf{400 tractors}, achieving real-time inertial navigation at \textbf{100 Hz} across diverse operational conditions}
      \resumeItemListEnd
      
      \vspace{2pt}
      \textit{\small \textbf{Visual Localization}}
      \resumeItemListStart
        \resumeItem{Developed a \textbf{visual relocalization} feature for the current \textbf{SLAM} system, improving performance and reliability in \textbf{GPS denied} indoor environments}
        \resumeItem{Implemented a \textbf{visual localization} module, utilizing pre-saved \textbf{keyframes}/\textbf{fiduciary markers} and integrated the results into a \textbf{pose graph optimizer} to improve the trajectory accuracy to enhance the system's ability to relocalize efficiently}
      \resumeItemListEnd
      
      \vspace{2pt}
      \textit{\small \textbf{Go To Planner}}
      \resumeItemListStart
        \resumeItem{Developed a novel \textbf{map representation} which segments the navigation space into map elements like \textbf{Roads}, \textbf{Intersections} and \textbf{Free space}}
        \resumeItem{Developed a custom \textbf{RRT\textsuperscript{*}} based \textbf{local motion planner} which generates a set of navigable paths for a given map element}
        \resumeItem{Demonstrated a \textbf{Proof of concept (POC)} which can generate shortest feasible global path within \textbf{10 msec} compared to the existing approach which takes \textbf{5 sec} approximately}
      \resumeItemListEnd
      
      \vspace{2pt}
      \textit{\small \textbf{Localization Stack Maintenance}}
      \resumeItemListStart
        \resumeItem{\textbf{Depth Map compression:} Implemented a \textbf{lossless depth map compression} module based on \textbf{Fast Lossless Depth Image Compression}, accelerating depth map compression by \textbf{8 times}}
        \resumeItem{\textbf{Wheel radius Calibration:} Developed a \textbf{real-time tire radius calibration} module using wheel speeds to improve vehicle localization accuracy}
        \resumeItem{\textbf{Complementary filter:} Implemented a \textbf{complementary filter} based on \textbf{multi-sensor fusion} from the \textbf{Inertial Measurement Unit (IMU)} measurements to determine the vehicle's orientation, providing an \textbf{accuracy of 0.8 degrees}}
      \resumeItemListEnd

  \resumeSubHeadingListEnd


%-----------PROJECTS-----------
\section{Projects}
    \resumeSubHeadingListStart
      \resumeProjectHeading
          {\textbf{Indoor Mobile Robot} $|$ \emph{ROS, C++, Python, OpenCV, LIDAR, AprilTags}}{Aug 2022 -- Apr 2023}
          
          \vspace{2pt}
          \textit{\small \textbf{Localization Stack}}
          \resumeItemListStart
            \resumeItem{Designed and implemented a \textbf{multi-sensor localization} stack for an indoor mobile robot, fusing data from \textbf{Wheel encoders}, \textbf{IMU}, \textbf{2D-LIDAR}, and \textbf{AprilTags} to achieve accuracy within \textbf{10 cm}}
            \resumeItem{Developed an efficient algorithm to extract \textbf{line features} from \textbf{2D LIDAR} data and associate them with map elements such as \textbf{corners} and \textbf{walls}}
            \resumeItem{Implemented a \textbf{camera pose estimation} module based on \textbf{AprilTags} strategically positioned in the map for accurate relative pose calculation}
            \resumeItem{Designed and implemented an optimized \textbf{Extended Kalman Filter (EKF)} for real-time \textbf{sensor fusion}, enhancing position update efficiency through \textbf{Numba}-based optimization}
          \resumeItemListEnd
          
          \vspace{2pt}
          \textit{\small \textbf{Planning Stack}}
          \resumeItemListStart
            \resumeItem{Developed and implemented a comprehensive \textbf{navigation system} for an indoor mobile robot, incorporating a \textbf{global path planner} and a \textbf{local motion planner} with \textbf{obstacle avoidance}}
            \resumeItem{Implemented an \textbf{A\textsuperscript{*}}-based \textbf{global planning} module using an \textbf{occupancy grid map} along with a \textbf{trajectory sampling-based local motion planner} with \textbf{obstacle avoidance} for real-time navigation}
          \resumeItemListEnd
          
      % \resumeProjectHeading
      %     {\textbf{VLA-Based Robot Control} $|$ \emph{Python, PyTorch, ROS}}{Month Year -- Month Year}
      %     \resumeItemListStart
      %       \resumeItem{Implemented \textbf{Vision-Language-Action (VLA)} model for \textbf{robot control}, enabling natural language-guided manipulation tasks}
      %       \resumeItem{Integrated \textbf{pre-trained vision-language models} with robot action prediction networks for end-to-end visuomotor control}
      %       \resumeItem{Developed data collection and training pipeline for fine-tuning VLA models on custom manipulation tasks}
      %     \resumeItemListEnd
          
      % \resumeProjectHeading
      %     {\textbf{RRT\textsuperscript{*} based Path Planning} $|$ \emph{C++, Python, Pybind11}}{Oct 2024}
      %     \resumeItemListStart
      %       \resumeItem{Built an \textbf{RRT\textsuperscript{*}} based \textbf{motion planning} framework in \textbf{C++} from scratch, focusing on efficient \textbf{global path planning}}
      %       \resumeItem{Smoothened the final path using a combination of \textbf{Artificial potential fields} and \textbf{elastic functions} for smoother and more efficient trajectories}
      %       \resumeItem{Designed software architecture for quick prototyping and testing by using \textbf{\href{https://pybind11.readthedocs.io/}{\color{blue}Pybind11}} to connect the planner and simulator. Utilized \textbf{\href{https://github.com/features/actions}{\color{blue}GitHub Actions}}, \textbf{Git hooks} and \textbf{unit tests} for continuous code quality checks and standard enforcements}
      %     \resumeItemListEnd
    \resumeSubHeadingListEnd

%
%-----------PROGRAMMING SKILLS-----------
\section{Technical Skills}
 \begin{itemize}[leftmargin=0.15in, label={}]
    \small{\item{
     \textbf{Languages}{: C++, Python, C, MATLAB} \\
     \textbf{Frameworks}{: ROS, PyTorch, OpenCV, Ceres, Sophus, Eigen, GTSAM, FAISS, Isaac Sim, Isaac Lab} \\
     \textbf{Developer Tools}{: Git, Docker, gtest, unittest} \\
     \textbf{Hardware}{: NVIDIA Jetson Orin, Stereo Cameras, LiDAR, IMU, Wheel Encoders, AprilTags} \\
     \textbf{Specializations}{: Visual Inertial SLAM, Sensor Fusion, Motion Planning, Deep Learning, Reinforcement Learning, Quadruped Locomotion}
    }}
 \end{itemize}

%-----------INTERESTS-----------
\section{Personal Interests}
 \begin{itemize}[leftmargin=0.15in, label={}]
    \small{\item{
     Reading books, Music, Writing, Chess
    }}
 \end{itemize}


%-------------------------------------------
\end{document}